\section{Numerical Experiments}
\label{sec:numerics}

The simulations distinguish three practical requirements for emulation.
The window must contain enough history, the fitted model must use that
information effectively, and the measurements must be sufficiently precise.
We examine these requirements in sequence.

\subsection{Models and Evaluation}

Our main experiments compare a fixed reset-QELM with two-, three-, and
four-qubit target QRCs. The targets have random $XX$ couplings, transverse
fields, input-dependent drives, and local loss $\gamma_T$, where $T$ denotes
the target QRC. Inputs are sampled uniformly from a 33-point alphabet on
$[0,1]$. Target outputs are single-site Pauli expectations and same-axis
two-site correlations.

The six-qubit reset-QELM encodes $d\leq6$ recent inputs, evolves under
independent random $XX$ couplings, and records Pauli features at nested times.
The original sweeps vary window length $d$, feature budget $B\in\{18,36,72\}$ counting
observable--checkpoint expectations, and circuit-shot budget per input $Q$.
These experiments evaluate a fixed feature architecture. The growing
universal family in Section~\ref{sec:criterion} establishes a separate
existence result.

Each independent realization block, called a root, fixes the target
parameters, input streams, and feature map. The confirmatory two- and
three-qubit cohort uses 12 roots. Dense maps use disjoint seeds, and a
prospectively specified four-qubit follow-up uses 12 new roots. A separate
prospective extension uses five-qubit targets, 12 new roots, and a ten-qubit
QELM with $d\in\{2,4,6,8,10\}$ and $B\in\{30,60,120\}$.
This extension increases both the available context and the emulator size.
Cohorts are analyzed separately.

Training, validation, test, and diagnostic streams are separate.
Standardization uses training data only, and validation selects the ridge
readout. Appendix~\ref{sec:methods} specifies the dynamics and fitting procedure.

To distinguish missing history from residual prediction error, we use a
held-out common-suffix panel. Different old histories are followed by
identical recent inputs. Mean squared variation around their mean output
measures the old-history contribution. The squared distance between that
mean and the emulator prediction measures the residual contribution.
These terms add exactly at the level of mean squared error
[Eq.~\eqref{eq:panel-decomposition}]. Residual emulator error can include
feature limitations, finite-data fitting, regularization, and finite-branch
estimation of the suffix mean. This sampled diagnostic does not directly
estimate the worst-case quantity $D_d$ or the spectral radius $r_{\rm vis}$.

Normalized root mean squared error (NRMSE) of $0.10$ corresponds to 10\%
of the target standard deviation. A setting meets the history or emulation
criterion when the corresponding error is at most $0.10$ in at least
10 of 12 roots. Confidence intervals use paired root-block bootstraps.

\begin{figure*}[t]
 \includegraphics[width=\textwidth]{fig04_long_context_n5.pdf}
 \caption{Five-qubit targets and longer context.
 Separate prospective cohort: 12 new roots and a fixed ten-qubit reset-QELM.
 \textbf{a,b} Common-suffix NRMSE versus window length at $\gamma_T=1$ and 4,
 comparing the history floor, $B=120$ QELM, and quadratic delay on logarithmic error axes.
 \textbf{c} Total QELM common-suffix NRMSE versus feature budget at fixed $d=10$
 on a linear error axis, using the same roots. Points show root medians; whiskers
 span the 10th--90th percentiles. Darker gray dashed lines mark 10\%.
 The larger emulator register is an additional resource, not a matched-size advantage.}
 \label{fig:long-context}
\end{figure*}

\subsection{History Requirements}

Longer windows and stronger local loss reduce the measured influence of
older inputs. Under weak loss, the missing past remains relevant throughout
the tested window range. Increasing loss shifts the minimum sufficient
window toward shorter contexts [Figure~\ref{fig:resource-diagnosis}(a)].

The dense maps cover 11 loss rates and windows from one to six, giving
198 size--loss--window settings across the two-, three-, and four-qubit
targets. Across these sizes, no tested window meets the history criterion
for $\gamma_T\leq1$. At $\gamma_T=4$, the minimum sufficient window is two
inputs. All rootwise local-loss decay fits have negative slopes in both
confirmatory sizes and in the independent four-qubit follow-up.

Longer context also improves the fitted QELM in the preregistered transition
settings. At $B=72$ and $\gamma_T\in\{2,4\}$, increasing $d$ from one to six
reduces held-out QELM NRMSE by a root median of $78.45\%$
(95\% CI $[76.64,80.80]\%$), with improvement in all 12 roots.
This measures the change in total emulator error, including both historical
and residual contributions.

The measured history-frontier pattern is robust to alternative success
thresholds. The primary $0.10$, 10-of-12 rule was preregistered.
A post hoc analysis uses tolerances $0.08$, $0.10$, and $0.12$, with 9, 10,
or 11 required roots. Individual crossing points change, but all 27 measured
history frontiers remain non-increasing with loss. Where no tested window
succeeds, the required context remains unresolved beyond the tested range.
These findings describe the tested systems and do not establish an
asymptotic size law.

The separate five-qubit extension shows what additional context can recover
[Figure~\ref{fig:long-context}(a,b)]. Increasing $d$ from two to ten reduces
history NRMSE by root medians of $93.55\%$ at $\gamma_T=1$ and $99.63\%$
at $\gamma_T=4$, with improvement in all 12 roots at each loss.
At $\gamma_T=1$, windows of six, eight, and ten inputs meet the history
criterion. At $\gamma_T=4$, every tested window meets it. The remaining
question is how effectively a fitted model can use this information.

\subsection{Feature Capacity and Classical Comparisons}

Sufficient history makes accurate emulation possible in principle, but
the tested feature model can still limit its accuracy. At $\gamma_T=4$,
a one-input window is history limited in every root even with $B=72$.
At $d=6$, history is sufficient, but the $B=18$ QELM usually fails while
replay succeeds. This is the operational representation-limit diagnosis
defined in Appendix~\ref{sec:methods}. With $B=36$, every root across the
two-, three-, and four-qubit targets reaches the tolerance.

Increasing feature capacity improves the residual error on this
history-sufficient setting. At $\gamma_T=4$ and $d=6$, increasing $B$ from
18 to 72 reduces common-suffix residual emulator NRMSE by a root median
of $67.26\%$ (95\% CI $[55.33,75.63]\%$).
The improvement occurs in all 12 roots in both confirmatory sizes
in the original cohort. All 24 confirmatory root-by-size
cases meet the emulation tolerance at $B=72$. The old-history term is unchanged.

A frozen follow-up extends the sweep to $B=144$ on 12 fresh roots per size
at $\gamma_T=4,d=6$, with fixed duration [Figure~\ref{fig:resource-diagnosis}(b)].
Exploratory $B=24,30$ reuse these data, retaining 18 final-time observables
and adding one or two per qubit at the midpoint with outcome-independent
Pauli-axis balance. Median error already declines before 36 features,
without isolating an approximation floor.

The five-qubit extension makes the distinction between history and fitted
accuracy particularly clear. At $\gamma_T=1$, no tested QELM budget meets
the total-error rule, even for the windows that meet the history criterion.
At $\gamma_T=4$, $B=60$ reaches the total-error tolerance in every root from
$d=4$ onward. At fixed $d=10$, most of the observed median-error reduction
occurs between 30 and 60 features, with little further decrease at 120
[Figure~\ref{fig:long-context}(c)]. These changes reflect the performance of
the fitted emulator, including feature approximation, estimation, and
regularization.

A longer window also need not improve fitted accuracy monotonically.
For each of the two-, three-, and four-qubit target sizes at $\gamma_T=2$
and $B=72$, $d=5$ meets the total-error rule while $d=6$ does not.
Window sufficiency and successful emulation must therefore be evaluated
separately.

Classical controls show how strongly the choice of model matters.
Linear delay, quadratic delay, and dimension-matched random-tanh models
use the same windows and data splits as the QELM. Quadratic delay is the
strongest classical control in all 24 fast-forgetting confirmatory cases
and outperforms the $B=72$ QELM in every root. In the five-qubit extension,
it also outperforms the $B=120$ QELM in 118 of 120 streaming root--condition
comparisons. These experiments demonstrate practical emulation limits
without establishing a quantum advantage.

\subsection{Measurement Budgets}

Additional features become useful only when measurements resolve them
accurately enough. We examine this requirement in a separate exploratory,
post-confirmatory study with 12 four-qubit targets at $\gamma_T=4$ and $d=6$,
where history is sufficient.

For each root, 16 repetitions independently sample training, validation,
and test features and repeat the complete fitting and selection procedure.
Validation selects both the ridge penalty and the feature budget.
Target outputs remain exact expectations, isolating the effect of finite
sampling in the emulator.

Median validation-selected test NRMSE falls from $0.231$ at 300 shots per
input to $0.0800$ at 19,200 shots [Figure~\ref{fig:resource-diagnosis}(c)].
The root median reduction is $65.20\%$ (95\% CI $[64.04,66.47]\%$),
with improvement in all 12 roots. Validation selects $B=36$ most often at
low shot budgets and $B=72$ at high budgets.

The useful feature budget therefore depends on the available measurement
budget. Once sufficient history is available, increasing feature capacity
and measurement precision together can improve practical emulation
\cite{hu2023sampling,ahmed2025finite}.
